\chapter{المراسلات الرسمية}
\label{ch:correspondence}

\section{ماذا ستنجز؟}

تجمع المراسلة الرسمية نص الكتاب والجهات المخاطبة ومسار العمل في سجل يمكن متابعته. تتعلم في هذه المرحلة تحديد الصفة التي تعمل باسمها والعثور على المراسلات التي تحتاج إلى انتباهك، والتمييز بين حالة الكتاب والإجراء المطلوب منك.

\begin{rcnote}[مثال: شركة الأفق]
يُعدّ عمر حداد طلب تنسيق زيارة فرع حلب، وتراجعه سارة الخطيب إذا أُسندت إليها المراجعة. يحدد كل منهما صفته قبل متابعة العمل؛ وهذا سيناريو تدريبي افتراضي، لا وصف لسجلات بيئة الاختبار.
\end{rcnote}

\begin{figure}[H]
\centering
\begin{tikzpicture}[node distance=0.5cm]
\node[stepBox] (a) {\RL{اختر الصفة}};
\node[stepBox,below=of a] (b) {\RL{حدّد العمل المطلوب}};
\node[stepBox,below=of b] (c) {\RL{افتح المراسلة وافحص حالتها}};
\draw[flowsTo] (a) -- (b);
\draw[flowsTo] (b) -- (c);
\end{tikzpicture}
\rclegend
\caption{بداية العمل بحسب المسؤولية}
\label{fig:correspondence-start}
\shotnote{تبدأ المتابعة بتحديد الصفة التي تعمل باسمها، ثم اختيار العمل الذي يخصها. افحص حالة المراسلة قبل اتخاذ أي إجراء؛ فاختيار مجلد أو فتح كتاب لا يعني اعتماده أو إنجازه.}
\end{figure}

\section{دليل صفحات البداية}
\label{sec:correspondence-pages}
ابدأ بالصفحات الثلاث بالترتيب. إعداد الكتاب واعتماده وأعمال الديوان والتهيئة مراحل لاحقة من المخطط؛ لا تحتاج إلى صلاحية الإدارة لتتعلم متابعة عملك اليومي.

\begin{table}[H]
\centering
\caption{من معرفة المسؤولية إلى العثور على العمل}
\label{tab:correspondence-pages-start}
\begin{tabularx}{\linewidth}{R{4.2cm} Y}
\rcHeaderRow
\thh{الصفحة} & \thh{ماذا تقدم هذه الصفحة؟}\\
الرئيسية ومسارك بحسب المسؤولية & تعرف أين يبدأ المُعِدّ والمعتمد والمستلم وموظف الديوان.\\
الصفة والصلاحيات & تختار المنصب وتفهم حدود ما يتيحه لك.\\
صندوق المراسلات ودورة العمل & تختار المجلد وتبحث وتقرأ الحالة والأولوية والسرية.\\
\end{tabularx}
\end{table}
